% GENERATED FILE. Edit canonical lessons, then run npm run generate:books.
% canonical-source-hash: fnv1a64:a4ebcf8f901797e0
% canonical-lessons: FR-C414-couchette, FR-C414-terminus, FR-C414-consigne, FR-C414-chariot, FR-C414-navette

\chapter{Le train de nuit}
\label{ch:fr-le-train-de-nuit}
\hlchaptermodality{414}

\begin{chapteropening}
\textbf{By the end of this chapter:} \emph{I can travel on a night train: the berth, the end of the line, the left-luggage office, the luggage trolley and the shuttle.}

Complete the last lesson of chapter 414: I can travel on a night train: the berth, the end of the line, the left-luggage office, the luggage trolley and the shuttle.
\end{chapteropening}

\section[la couchette]{la couchette — the berth, the couchette}

\label{lesson:FR-C414-couchette}



\begin{quote}
Before the new one: say the French for the carriage, then the French for to stop at.
\end{quote}

\subsection*{You'll want to know: la couchette}
\textbf{la couchette} — \textquotedblleft{}the berth, the couchette\textquotedblright{}.

\begin{quote}
Dans le train de nuit, j'ai dormi sur la couchette du haut.
\end{quote}

\begin{quote}
\emph{On the night train I slept on the top berth.}
\end{quote}

A small \emph{couche} 'bed'.

\begin{cousinweb}
\emph{couchette} is a small \emph{couche} 'bed', from \emph{coucher} 'to lay down', which is inherited from Latin \emph{collocāre} 'to put in place', built on \emph{locus} 'place' (Lewis \& Short).
\end{cousinweb}

\begin{grammarlens}[title={What you've built}]
One more word for telling, explaining and describing at B1.
\end{grammarlens}

\subsection*{Your turn}
\begin{itemize}
\raggedright
  \item \emph{Say it:} \emph{la couchette}
  \item \emph{Say it:} \emph{la couchette}, once more
  \item \emph{Say it:} a sentence of your own with \emph{la couchette}
  \item \emph{Recall:} say the French for the carriage, then the French for to stop at, then say \emph{la couchette} again
\end{itemize}

\begin{tcolorbox}[breakable,colback=teal!4,colframe=teal!35!black,title={Before you move on}]
What does la couchette mean? (\textquotedblleft{}The berth, the couchette\textquotedblright{}.) Say it once more.
\end{tcolorbox}
\section[le terminus]{le terminus — the end of the line, the terminus}

\label{lesson:FR-C414-terminus}



\begin{quote}
Before the new one: say the French for to stop at, then the French for the berth.
\end{quote}

\subsection*{You'll want to know: le terminus}
\textbf{le terminus} — \textquotedblleft{}the end of the line, the terminus\textquotedblright{}.

\begin{quote}
Le bus s'arrête au terminus et tout le monde descend.
\end{quote}

\begin{quote}
\emph{The bus stops at the end of the line and everyone gets off.}
\end{quote}

The final \emph{s} is pronounced.

\begin{cousinweb}
\emph{terminus} was borrowed from English \emph{terminus}, a learned word from Latin \emph{terminus} 'boundary, end' (EtymDB).
\end{cousinweb}

\begin{grammarlens}[title={What you've built}]
One more word for telling, explaining and describing at B1.
\end{grammarlens}

\subsection*{Your turn}
\begin{itemize}
\raggedright
  \item \emph{Say it:} \emph{le terminus}
  \item \emph{Say it:} \emph{le terminus}, once more
  \item \emph{Say it:} a sentence of your own with \emph{le terminus}
  \item \emph{Recall:} say the French for to stop at, then the French for the berth, then say \emph{le terminus} again
\end{itemize}

\begin{tcolorbox}[breakable,colback=teal!4,colframe=teal!35!black,title={Before you move on}]
What does le terminus mean? (\textquotedblleft{}The end of the line, the terminus\textquotedblright{}.) Say it once more.
\end{tcolorbox}
\section[la consigne]{la consigne — the left-luggage office; the instruction}

\label{lesson:FR-C414-consigne}



\begin{quote}
Before the new one: say the French for the berth, then the French for the end of the line.
\end{quote}

\subsection*{You'll want to know: la consigne}
\textbf{la consigne} — \textquotedblleft{}the left-luggage office; the instruction\textquotedblright{}.

\begin{quote}
Nous avons laissé nos sacs à la consigne.
\end{quote}

\begin{quote}
\emph{We left our bags at the left-luggage office.}
\end{quote}

\emph{Les consignes de sécurité} — the safety instructions.

\begin{cousinweb}
\emph{consigne} is built on \emph{consigner} 'to put in safe keeping', a learned borrowing from Latin \emph{consignāre} 'to seal', from \emph{signum} 'mark' (Lewis \& Short).
\end{cousinweb}

\begin{grammarlens}[title={What you've built}]
One more word for telling, explaining and describing at B1.
\end{grammarlens}

\subsection*{Your turn}
\begin{itemize}
\raggedright
  \item \emph{Say it:} \emph{la consigne}
  \item \emph{Say it:} \emph{la consigne}, once more
  \item \emph{Say it:} a sentence of your own with \emph{la consigne}
  \item \emph{Recall:} say the French for the berth, then the French for the end of the line, then say \emph{la consigne} again
\end{itemize}

\begin{tcolorbox}[breakable,colback=teal!4,colframe=teal!35!black,title={Before you move on}]
What does la consigne mean? (\textquotedblleft{}The left-luggage office; the instruction\textquotedblright{}.) Say it once more.
\end{tcolorbox}
\section[le chariot]{le chariot — the trolley (for luggage or shopping)}

\label{lesson:FR-C414-chariot}



\begin{quote}
Before the new one: say the French for the end of the line, then the French for the left-luggage office.
\end{quote}

\subsection*{You'll want to know: le chariot}
\textbf{le chariot} — \textquotedblleft{}the trolley (for luggage or shopping)\textquotedblright{}.

\begin{quote}
Il pousse un chariot plein de valises.
\end{quote}

\begin{quote}
\emph{He is pushing a trolley full of suitcases.}
\end{quote}

In a supermarket, \emph{le chariot} is the shopping trolley.

\begin{cousinweb}
\emph{chariot} is built on Old French \emph{char} 'cart', which is inherited from Latin \emph{carrus} 'wagon' (EtymDB).
\end{cousinweb}

\begin{grammarlens}[title={What you've built}]
One more word for telling, explaining and describing at B1.
\end{grammarlens}

\subsection*{Your turn}
\begin{itemize}
\raggedright
  \item \emph{Say it:} \emph{le chariot}
  \item \emph{Say it:} \emph{le chariot}, once more
  \item \emph{Say it:} a sentence of your own with \emph{le chariot}
  \item \emph{Recall:} say the French for the end of the line, then the French for the left-luggage office, then say \emph{le chariot} again
\end{itemize}

\begin{tcolorbox}[breakable,colback=teal!4,colframe=teal!35!black,title={Before you move on}]
What does le chariot mean? (\textquotedblleft{}The trolley (for luggage or shopping)\textquotedblright{}.) Say it once more.
\end{tcolorbox}
\section[la navette]{la navette — the shuttle (bus or train)}

\label{lesson:FR-C414-navette}



\begin{quote}
Before the new one: say the French for the left-luggage office, then the French for the trolley.
\end{quote}

\subsection*{You'll want to know: la navette}
\textbf{la navette} — \textquotedblleft{}the shuttle (bus or train)\textquotedblright{}.

\begin{quote}
Une navette relie le centre-ville à l'aéroport.
\end{quote}

\begin{quote}
\emph{A shuttle links the town centre with the airport.}
\end{quote}

\emph{Faire la navette} — to go back and forth.

\begin{cousinweb}
\emph{navette} is a small \emph{nef} 'ship', from Latin \emph{nāvis} 'ship' (Lewis \& Short): first the boat-shaped shuttle of a loom that runs back and forth.
\end{cousinweb}

\begin{grammarlens}[title={What you've built}]
One more word for telling, explaining and describing at B1.
\end{grammarlens}

\subsection*{Your turn}
\begin{itemize}
\raggedright
  \item \emph{Say it:} \emph{la navette}
  \item \emph{Say it:} \emph{la navette}, once more
  \item \emph{Say it:} a sentence of your own with \emph{la navette}
  \item \emph{Recall:} say the French for the left-luggage office, then the French for the trolley, then say \emph{la navette} again
\end{itemize}

\begin{tcolorbox}[breakable,colback=teal!4,colframe=teal!35!black,title={Before you move on}]
What does la navette mean? (\textquotedblleft{}The shuttle (bus or train)\textquotedblright{}.) Say it once more.
\end{tcolorbox}
